\section{Convex-body criteria for varying references}
\label{sec:convex-body-reference-stability}
\index{Reference stability!convex-body criterion}

The closed mass formula gives a geometric convergence criterion when a convex body describes the singularities. A common translated simplex ensures that the limiting reference remains solvable. Hausdorff distance then controls the functional uniformly over all continuous tests, even when the ordinary types vary.

Let $n\geq2$, and let $\omega_{\mathrm{FS}}$ be the Fubini--Study form on $\mathbb P^n$, normalized by $\int\omega_{\mathrm{FS}}^n=1$. On the torus $T=\{Z_0\cdots Z_n\ne0\}$ put
\[
 x_i=\log|Z_i/Z_0|^2,\qquad
 L(x)=\log\bigl(1+\textstyle\sum_i\mathrm{e}^{x_i}\bigr),\qquad
 \Sigma=\{p_i\geq0:\textstyle\sum_i p_i\leq1\}.
\]
Thus $\omega_{\mathrm{FS}}=\ddc L$ in this chart. Define the continuous simplex entropy
\[
 \ell(p)=\sum_i p_i\log p_i+
       (1-\textstyle\sum_i p_i)\log(1-\textstyle\sum_i p_i),
 \qquad 0\log0=0.
\]
For a convex body $K\subset\Sigma$, set
\[
 h_K(x)=\sup_{p\in K}p\cdot x,\qquad
 g_K(x)=\sup_{p\in K}\{p\cdot x-\ell(p)\}.
\]

\subsection{Canonical models and their mass}

\begin{lemma}\label{lem:convex-body-model}
The functions $h_K-L$ and $g_K-L$ extend canonically to $\omega_{\mathrm{FS}}$-plurisubharmonic functions $\psi_K,\Phi_K$ on $\mathbb P^n$. They satisfy
\[
 \begin{aligned}
  &\psi_K\leq\Phi_K\leq\psi_K+\log(n+1),\qquad\sup\Phi_K=0,\\
  &P_{\omega_{\mathrm{FS}}}[\psi_K]=\Phi_K
       =P_{\omega_{\mathrm{FS}}}[\Phi_K].
 \end{aligned}
\]
If $p_*+\delta\Sigma\subset K$ for $0<\delta<1$, then $P_{(1-\delta)\omega_{\mathrm{FS}}}(\Phi_K)$ is genuine.
\end{lemma}
\begin{proof}
Convex conjugation gives $L^*=\ell$ on $\Sigma$ and $+\infty$ outside it. Since $-\log(n+1)\leq\ell\leq0$,
\[
 h_K\leq g_K\leq h_K+\log(n+1),\qquad
 g_K\leq L,\qquad h_K\leq h_\Sigma\leq L.
\]
A finite convex function of the logarithmic moduli is plurisubharmonic: for a smooth convex function the local complex Hessian is positive, and convolution gives the general case. In the chart $Z_j\ne0$, add the local Fubini--Study potential. The resulting functions are bounded above by the smooth chart potential because $h_K,g_K\leq L$. \Cref{lem:qpsh-removable-analytic} therefore gives the canonical extensions and the stated bounds. If $p\in\operatorname{int}K$, it lies in $\operatorname{int}\Sigma$. At $x_i=\log(p_i/(1-\sum p_i))$, conjugacy gives $g_K(x)=L(x)$, proving $\sup\Phi_K=0$.

For $C\geq\log(n+1)$, let $v_C=P_{\omega_{\mathrm{FS}}}(\psi_K+C,0)$. The actual competitor $\Phi_K$ gives $\Phi_K\leq v_C\leq0$. The obstacle is upper semicontinuous, so $v_C\leq\psi_K+C$ everywhere. Rotations preserve the obstacle; hence $v_C$ is torus invariant. Its finite local potential $f=v_C+L$ is convex and satisfies $f\leq L$ and $f\leq h_K+C$. Separation shows $\operatorname{dom}f^*\subset K$, while the first inequality gives $f^*\geq\ell$. Thus
\[
 f(x)=\sup_{p\in\operatorname{dom}f^*}\{p\cdot x-f^*(p)\}
       \leq g_K(x).
\]
Canonical equality off the coordinate boundary gives $v_C=\Phi_K$ globally. The finite rooftop families for bounded-equivalent obstacles interlace, proving both ordinary model identities.

For the final assertion put $\varepsilon=1-\delta$, $b=\sum_i p_{*,i}$ and $\chi=p_*\cdot x-\varepsilon L$. The simplex inclusion implies $p_{*,i}\geq0$ and $b\leq\varepsilon$. Its canonical extension satisfies
\[
 \varepsilon\omega_{\mathrm{FS}}+\ddc\chi
 =\sum_i p_{*,i}[Z_i=0]+(\varepsilon-b)[Z_0=0]\geq0.
\]
Since $h_K\geq p_*\cdot x+\delta h_\Sigma$ and $L-h_\Sigma\leq\log(n+1)$, one has $\chi-\delta\log(n+1)\leq\Phi_K$. This genuine smaller-background competitor proves the claim.
\end{proof}

Take $X=\mathbb P^n$ with $\pi=\mathrm{id}$, or a fixed point blowup $\pi:X=\operatorname{Bl}_q\mathbb P^n\to\mathbb P^n$, and put $\Theta=\pi^*\omega_{\mathrm{FS}}$.

\begin{lemma}\label{lem:convex-body-mass}\index{Convex body!model mass}
The manifold $X$ is K\"ahler and $\Theta$ is smooth semipositive and big. For $r>0$, the reference $r\pi^*\Phi_K$ is an ordinary model for $r\Theta$, has small unbounded locus, and satisfies
\begin{equation}\label{eq:convex-body-mass}
 \int_X\mathrm{MA}_{r\Theta}(r\pi^*\Phi_K)
       =r^n n!\operatorname{Vol}_n(K).
\end{equation}
The smaller-background witness in \cref{lem:convex-body-model} also scales and pulls back.
\end{lemma}
\begin{proof}
On $\mathbb P^n$ the model assertion is \cref{lem:convex-body-model}. In the blowup case, a genuine $\pi^*\omega_{\mathrm{FS}}$-plurisubharmonic potential descends under the point blowup. Off its exceptional divisor, add a local downstairs background potential. The resulting plurisubharmonic function is upper bounded near the center, since the upstairs potential is upper bounded near the compact fiber. Removable extension gives a downstairs potential whose pullback agrees with the original off the divisor, hence everywhere canonically. The inverse is pullback. This bijection preserves inequalities and envelopes, and positive scaling preserves them as well.

On projective space $\Theta=\omega_{\mathrm{FS}}$ is K\"ahler. For bigness in the blowup case, choose unitary homogeneous coordinates with $q=[1:0:\cdots:0]$, set
\[
 \psi_q=\log\frac{\sum_{i=1}^n|Z_i|^2}{\sum_{i=0}^n|Z_i|^2},
\]
and let $Q:X\to\mathbb P^{n-1}$ be the resolved direction map. For $0<t<1$,
\[
 \Theta+\ddc(t\pi^*\psi_q)
 =(1-t)\pi^*\omega_{\mathrm{FS}}+tQ^*\omega_{\mathrm{FS},n-1}+t[E_q].
\]
The smooth sum is K\"ahler, since $(\pi,Q)$ is the incidence graph embedding and both coefficients are positive. This proves the assertions about the manifold and class.

For a smooth convex $g$, the coordinates $z_i=\mathrm{e}^{x_i/2+\mathrm{i}\vartheta_i}$ give
\[
 \mathrm{d}^{\mathrm c}g=\frac1{2\pi}\sum_i g_i\,\mathrm{d}\vartheta_i,
 \qquad
 \int_T(\ddc g)^n=n!\int_{\mathbb R^n}\det(D^2g)\,\mathrm{d}x.
\]
Convolution and local uniform continuity of bounded plurisubharmonic products identify the same measure for finite convex $g$: its real factor is the Alexandrov measure $M_g$, which assigns to a Borel set the Euclidean volume of its subgradient image. For $g_K$, separation and $g_K\leq h_K+\log(n+1)$ put every subgradient in $K$. Conversely, if $p\in\operatorname{int}K$ and $p+cB\subset K$, then $g_K(x)-p\cdot x\geq c|x|$; it attains a minimum, so $p$ is a subgradient somewhere. Consequently
\[
 \operatorname{int}K\subset\partial g_K(\mathbb R^n)\subset K,
 \qquad M_{g_K}(\mathbb R^n)=\operatorname{Vol}_n(K).
\]
The boundary of a convex body has zero volume. The analytic complement of $T$ receives no non-pluripolar mass, so this is the full mass. Locality off the exceptional divisor proves the pullback formula; that divisor and its image carry no non-pluripolar mass. Homogeneity gives the factor $r^n$. The reference is locally bounded off the analytic coordinate boundary and its inverse image, hence has small unbounded locus. The explicit witness pulls back and scales as a genuine potential.
\end{proof}

\Needspace{16\baselineskip}
\subsection{Hausdorff control of the functional}

\begin{theorem}\label{thm:convex-body-reference-stability}
Let convex bodies $P_j,P\subset\Sigma$ contain the same translated simplex $p_*+\delta\Sigma$, with $0<\delta<1$, and suppose $h_j=d_H(P_j,P)\to0$. Let $r_j>0$ tend to $r>0$. On the fixed $X$ above set
\[
 (\theta_j,\varphi_j)=(r_j\Theta,r_j\pi^*\Phi_{P_j}),
 \qquad (\theta,\varphi)=(r\Theta,r\pi^*\Phi_P).
\]
Each datum, including the limit, has an ordinary-type solution for every bounded real Borel $F$. Its mass is $m_j=r_j^n n!\operatorname{Vol}_n(P_j)$ or $m=r^n n!\operatorname{Vol}_n(P)>0$. For a fixed smooth positive volume $\mathrm{d}V$ on $X$, put $Z_F=\int_X\mathrm{e}^F\,\mathrm{d}V$. Then
\[
 B_j(F)=\log m_j-\log Z_F,\qquad B(F)=\log m-\log Z_F.
\]
Writing $s=\delta/(n+\sqrt n)$, one has
\begin{equation}\label{eq:convex-body-functional-bound}
 \sup_{G\in C(X,\mathbb R)}|B_j(G)-B(G)|
 \leq n|\log(r_j/r)|+n\log(1+h_j/s).
\end{equation}
Moreover $n!\operatorname{Vol}_n(P_j),n!\operatorname{Vol}_n(P)\geq\delta^n$, the backgrounds converge smoothly, and the references converge in $L^1(X)$ and locally uniformly on $\pi^{-1}(T)$. For continuous real $F$, every supporting set equals $\{\nu_F\}$, where $\nu_F=\mathrm{e}^F\mathrm{d}V/Z_F$. If continuous $F_j\to F$ uniformly and $\nu_j\in\mathcal A_j(F_j)$, then
\begin{equation}\label{eq:convex-body-support-bound}
 \int_X\left|\frac{\mathrm{d}\nu_j}{\mathrm{d}V}
          -\frac{\mathrm{d}\nu_F}{\mathrm{d}V}\right|\mathrm{d}V
 \leq\mathrm{e}^{2\|F_j-F\|_\infty}-1.
\end{equation}
\end{theorem}
\begin{proof}
The preceding lemmas give the actual limiting model, its smaller-background witness and its positive mass. \Cref{thm:normalization-mass-quotient} therefore gives ordinary-type solvability and the scalar and supporting formulas at every datum.

The common simplex has an inscribed ball $B(c,s)$. Hausdorff distance gives $P_j\subset P+h_jB$, while $sB\subset P-c$ and convexity give
\[
 P+h_jB\subset c+(1+h_j/s)(P-c).
\]
Reversing $P,P_j$ bounds both volume ratios by $(1+h_j/s)^n$. Subtract the scalar formulas and include the factor $r_j^n/r^n$ to obtain \eqref{eq:convex-body-functional-bound}. The simplex volume is $\delta^n/n!$.

Let $\omega_\ell$ be the uniform-continuity modulus of $\ell$ on $\Sigma$. Matching entropy maximizers by Hausdorff distance gives
\[
 |g_{P_j}(x)-g_P(x)|\leq h_j|x|+\omega_\ell(h_j).
\]
This proves local uniform convergence on $T$. The common lower bound $\chi-\delta\log(n+1)$ from \cref{lem:convex-body-model} is integrable on $X$, including after pullback. Boundedness of $r_j$ and dominated convergence off the analytic boundary give $L^1$ convergence. Finally the normalized exponential densities have ratio in $[\mathrm{e}^{-2\eta},\mathrm{e}^{2\eta}]$, where $\eta=\|F_j-F\|_\infty$. This proves \eqref{eq:convex-body-support-bound} and convergence in total variation.
\end{proof}

\subsection{A family with nonanalytic isolated poles}
\label{sec:curved-nonanalytic-references}
\index{Singularity type!nonanalytic isolated pole}

The common-simplex hypothesis allows ordinary singularities that have no analytic representative up to bounded error. For $0<a\leq1$, define
\[
 P_a=\{p\in\Sigma:\textstyle\sum_i\sqrt{p_i}\geq\sqrt a\},
 \qquad
 \kappa_n=\frac{2^n n!}{(2n)!}=\frac1{(2n-1)!!}.
\]

\begin{proposition}\label{prop:curved-nonanalytic-reference}
The model $\Phi_{P_a}$ has exactly the isolated pole $p=[1:0:\cdots:0]$ and is not bounded-equivalent to an analytic singularity. Its generic Lelong number along every divisor of $\mathbb P^n$ is zero, and
\begin{equation}\label{eq:curved-nonanalytic-mass}
 \int_{\mathbb P^n}\mathrm{MA}_{\omega_{\mathrm{FS}}}(\Phi_{P_a})
 =1-\kappa_n a^n>0.
\end{equation}
The same pole, nonanalyticity and zero generic-divisor assertions hold on a fixed blowup at $q\ne p$.
\end{proposition}
\begin{proof}
The square-root sum is concave, so $P_a$ is convex. Set
\[
 b=\frac{n+1}{2n},\qquad p_{*,i}=\frac bn,
 \qquad\delta=\frac{n-1}{4n}.
\]
Then $b+\delta<1$ and $\sum_i\sqrt{p_{*,i}}=\sqrt{(n+1)/2}>1$. Adding a point of $\delta\Sigma$ preserves that inequality, giving a common translated simplex in every $P_a$.

All vertices $e_i$ belong to $P_a$, and therefore
\[
 \psi_{P_a}\geq\log\frac{\max_{1\leq i\leq n}|Z_i|^2}
                                  {\sum_{i=0}^n|Z_i|^2}.
\]
The reference is locally bounded away from $p$. Cauchy--Schwarz gives $\sum p_i\geq a/n$ on $P_a$. When all $x_i<0$, this yields $h_{P_a}(x)\leq(a/n)\max_i x_i\to-\infty$ as $z\to p$. \Cref{lem:convex-body-model} gives the claimed pole. Local boundedness outside one point proves the generic divisorial assertion; pullback at a distinct point introduces no pole.

For positive $\lambda_i$ take $|z_i|^2=\mathrm{e}^{-R\lambda_i}$. Writing $y_i=\sqrt{p_i}$, Cauchy--Schwarz gives
\begin{equation}\label{eq:curved-directional-coefficient}
 \min_{p\in P_a}\sum_i\lambda_i p_i
 =\frac a{\sum_i\lambda_i^{-1}}=:d_a(\lambda),\qquad
 y_i=\frac{\sqrt a}{\lambda_i\sum_k\lambda_k^{-1}}.
\end{equation}
The minimizer has $\sum p_i\leq a\leq1$. Since $L(-R\lambda)\to0$ and $\Phi_{P_a}-\psi_{P_a}$ is bounded, $-\Phi_{P_a}/R\to d_a(\lambda)$ along these directions.

Suppose the type were $c\log\sum_k|f_k|^2+O(1)$, with $c>0$ and holomorphic germs $f_k=\sum_\alpha c_{k,\alpha}z^\alpha$. Exponentiating the bounded-equivalence inequalities and averaging phases gives the positive series
\[
 \sum_{k,\alpha}|c_{k,\alpha}|^2\mathrm{e}^{-R\alpha\cdot\lambda},
\]
whose exponential rate must be $d_a(\lambda)/c$. That rate is the minimum of the monomial exponents: a minimizing term gives one bound, and convergence at a fixed $R_0$ bounds the series above by a constant times $\mathrm{e}^{-(R-R_0)\min\alpha\cdot\lambda}$. Near any strictly positive $\lambda$, only finitely many monomials can minimize, because the coordinates have a common positive lower bound and one nonzero coefficient bounds the active degree. The rate is therefore locally piecewise linear. But
\[
 d_a(s,1,\ldots,1)=\frac{as}{1+(n-1)s},\qquad
 \frac{\mathrm{d}^2}{\mathrm{d}s^2}d_a(s,1,\ldots,1)
 =-\frac{2a(n-1)}{(1+(n-1)s)^3}<0.
\]
This contradiction proves nonanalyticity. The assertion persists on a blowup away from $p$, where a neighborhood of the pole is unchanged.

The removed corner $C_a=\{p_i\geq0:\sum_i\sqrt{p_i}<\sqrt a\}=aC_1$ lies in $\Sigma$. Substitution $p_i=y_i^2$ gives
\[
 \operatorname{Vol}_n(C_1)
 =2^n\int_{y_i\geq0,\,\sum y_i<1}\prod_i y_i\,\mathrm{d}y
 =\frac{2^n}{(2n)!}.
\]
The last equality follows by iterated monomial integration over a simplex. Subtract the corner from $\Sigma$ and apply \eqref{eq:convex-body-mass} to obtain \eqref{eq:curved-nonanalytic-mass}. Since $n\geq2$, one has $\kappa_n<1$.
\end{proof}

\begin{corollary}\label{cor:curved-reference-stability}
Put $a_j=1-1/(j+2)$ and $r_j=1+1/(j+2)$, $j\geq1$. On $\mathbb P^n$ or its fixed blowup at $q\ne p$, the models $r_j\pi^*\Phi_{P_{a_j}}$ satisfy \cref{thm:convex-body-reference-stability}, with limit $\pi^*\Phi_{P_1}$. Their ordinary types and masses vary, and
\[
 \sup_G|B_j(G)-B(G)|
 \leq n|\log r_j|+\frac{n\kappa_n}{1-\kappa_n}|1-a_j|\longrightarrow0.
\]
The limiting class is solvable and has mass $1-\kappa_n$, equal to $2/3$ for $n=2$.
\end{corollary}
\begin{proof}
If $p\in P_{a_j}\setminus P_1$, then $u=(\sum\sqrt{p_i})^2\in[a_j,1)$ and $p/u\in P_1$. Since $|p|\leq\sum p_i\leq u$, one has $|p/u-p|\leq1-u\leq1-a_j$. Thus $d_H(P_{a_j},P_1)\leq1-a_j$. The mass formula and the derivative bound $|\mathrm{d}\log(1-\kappa_n a^n)/\mathrm{d}a|\leq n\kappa_n/(1-\kappa_n)$ give the displayed estimate. The directional coefficient is scaled by $r_ja_j=1-(j+2)^{-2}$, which strictly varies, while $r_j^n(1-\kappa_n a_j^n)$ strictly decreases to $1-\kappa_n>0$.
\end{proof}


\subsection{Positive mass with a degenerate smooth residual}
\label{sec:degenerate-projective-residuals}
\index{Reference stability!degenerate residual}

A second family has a logarithmic pole and a smooth resolved residual that degenerates along a different exceptional divisor. Fix $p=[1:0:\cdots:0]$ and $q\ne p$ on $\mathbb P^n$. On $X=\operatorname{Bl}_q\mathbb P^n$ keep $\Theta=\pi^*\omega_{\mathrm{FS}}$. For $0<t<1$, put
\[
 \psi_t=t\log\frac{\sum_{i=1}^n|Z_i|^2}{\sum_{i=0}^n|Z_i|^2},
 \qquad K_t=\{p\in\Sigma:\textstyle\sum_i p_i\geq t\}.
\]
The local current potential of $\psi_t$ is $(1-t)L+t\log\sum_i\mathrm{e}^{x_i}$. Its bounded-distance support function is
\[
 (1-t)h_\Sigma+t\max_i x_i=h_{K_t},\qquad
 K_t=(1-t)\Sigma+t\operatorname{conv}(e_1,\ldots,e_n).
\]
Thus $\psi_t\sim\psi_{K_t}\sim\Phi_{K_t}$, and $n!\operatorname{Vol}_n(K_t)=1-t^n$.

\begin{proposition}\label{prop:degenerate-projective-residual}
Let $\tau:Z\to X$ resolve the pole by blowing up the point above $p$. There is a smooth semipositive residual $\alpha_t$ such that
\begin{equation}\label{eq:degenerate-projective-residual}
 (\pi\circ\tau)^*\omega_{\mathrm{FS}}
       +\ddc(\psi_t\circ\pi\circ\tau)=\alpha_t+t[E_p].
\end{equation}
It is positive away from $E_q$, has rank exactly one at every point of $E_q$, and has top integral $1-t^n>0$. Every ordinary-type solution associated with $\pi^*\Phi_{K_t}$ lifts, after subtracting the prescribed divisor pole, to a bounded canonical $\alpha_t$-plurisubharmonic remainder.
\end{proposition}
\begin{proof}
On $Y=\operatorname{Bl}_p\mathbb P^n$, let $\pi_p$ be the blowdown and $Q_p:Y\to\mathbb P^{n-1}$ the resolved direction map. The incidence graph gives the K\"ahler form
\[
 \beta_t=(1-t)\pi_p^*\omega_{\mathrm{FS}}
             +tQ_p^*\omega_{\mathrm{FS},n-1}.
\]
The disjoint centers identify $Z$ with the blowup $\sigma:Z\to Y$ at the point above $q$. Put $\alpha_t=\sigma^*\beta_t$. Poincar\'e--Lelong gives \eqref{eq:degenerate-projective-residual}. There is no $E_q$ coefficient, since $\psi_t$ is locally bounded at $q$.

In a blowup chart $\sigma$ is $(s,w_2,\ldots,w_n)\mapsto(s,sw_2,\ldots,sw_n)$. At $s=0$ its differential has rank one and kills the $n-1$ tangent directions of $E_q$. Positivity of $\beta_t$ gives the exact rank assertion; off $E_q$ the map is biholomorphic.

Write $H=\pi_p^*[\omega_{\mathrm{FS}}]$ and $E=[E_p]$. The current identity gives $[\beta_t]=H-tE$. Since $H|_{E_p}=0$, all mixed top intersections vanish. The normal line of $E_p$ is $\mathcal O_{\mathbb P^{n-1}}(-1)$, so $\int E^n=(-1)^{n-1}$. Therefore
\[
 \int_Z\alpha_t^n=\int_Y(H-tE)^n=1-t^n.
\]
The first equality uses the degree-one blowup. An ordinary-type solution differs boundedly from the pulled-back prototype off $E_p$. After adding a local smooth residual potential, its bounded remainder is plurisubharmonic there and extends canonically across $E_p$ by \cref{lem:qpsh-removable-analytic}. This proves the last assertion without inserting a divisor value into an unbounded product.
\end{proof}

\begin{corollary}\label{cor:degenerate-residual-reference-stability}
Suppose $t_j\to t$ in $[t_{\min},t_{\max}]\subset(0,1)$, and let $r_j>0$ tend to $r>0$. The models $r_j\pi^*\Phi_{K_{t_j}}$ satisfy \cref{thm:convex-body-reference-stability}, including ordinary-type solvability at the limit, and
\[
 \sup_G|B_j(G)-B(G)|
 \leq n|\log(r_j/r)|+
       \frac{nt_{\max}^{n-1}}{1-t_{\max}^n}|t_j-t|.
\]
All continuous-test supports are the normalized exponential probabilities from that theorem. The zero-mass endpoint $t=1$ admits no positive-density equation in the prototype's ordinary type.
\end{corollary}
\begin{proof}
Choose $p_{*,i}=t_{\max}/n$ and $\delta=(1-t_{\max})/2$. Their translated simplex lies in every $K_t$ in the interval. For $s\geq t$, move a point of $K_t$ toward a vertex until its coordinate sum reaches $s$. This gives
\[
 d_H(K_s,K_t)\leq\frac{\sqrt2}{1-t_{\max}}|s-t|.
\]
Apply the stability theorem and use the derivative bound for $\log(1-t^n)$ to obtain the sharper displayed functional estimate.

At $t=1$, the resolved residual is pulled back from $\mathbb P^{n-1}$, so its top integral is zero. Every ordinary-type potential has a bounded residual-plurisubharmonic remainder. Bounded closed-product Stokes gives the same zero top mass for that remainder; positivity forces its top measure to vanish. Thus it cannot solve $c\mathrm{e}^F\mathrm{d}V$ with bounded $F$ and $c>0$. In contrast, the endpoint $a=1$ in \cref{cor:curved-reference-stability} has positive mass and remains solvable.
\end{proof}

Scaling by $r>0$ multiplies the resolved residual and the prescribed divisor coefficient by $r$. Solvability is established on $X$ by the model criterion, without treating the resolved residual as strictly positive. These criteria use closed projective backgrounds and fixed modifications. Convergence for general nonclosed Hermitian backgrounds, moving modifications or unrestricted singularity degeneration remains open. The supporting probabilities here identify the limiting functional; no convergence of solution potentials or rough weighted adjoint products is inferred.

